% why-self-modelers.tex — Section 3: Why Self-Modelers?

\section{Why Self-Modelers?}
\label{sec:why-self-modelers}

\subsection{Experiential measure}
\label{subsec:experiential-measure}

Structure space $\mathcal{S}$ contains every consistent mathematical
structure.  Most have no observers.  The question ``why do observers
experience this specific physics?'' requires a way to identify which
structures contain observers and how much experiential weight they
carry.

\begin{postulate}[Self-Modeling]
\label{post:self-modeling}
Experiential weight attaches to self-modeling structures: composite
Markov processes $(B, M)$ in which a body component $B$ is tracked by
a model component $M$ through a factorizable transition kernel.
\end{postulate}

Paper~1~\citep{ehrlich2026measure} defines the experiential density
\begin{equation}
\label{eq:rho}
\rho(p) = I_p(B;\,M)\,\bigl(1 - I_p(B;\,M) / H_p(B)\bigr),
\end{equation}
where $I_p(B;\,M)$ is the mutual information between body and model at
state $p$, and $H_p(B)$ is the body's entropy.  This density has three
qualitative features:
\begin{enumerate}[nosep]
\item $\rho = 0$ when $I = 0$: no model, no experience.
\item $\rho = 0$ when $I = H$: perfect self-knowledge, nothing left
  to experience (the model is trivially the system).
\item $\rho > 0$ at intermediate fidelity: imperfect but active
  self-modeling.
\end{enumerate}
Rocks, thermostats, and unstructured noise score $\rho = 0$.  Cells
score near zero.  Brains with rich self-models score high.  Boltzmann
brains --- brief thermal fluctuations --- score near zero when
integrated over duration, because their self-modeling is fleeting
(Paper~2 proves this rigorously via metastability
theory~\citep{ehrlich2026bb}).

\subsection{The $\rho$ bootstrap}
\label{subsec:rho-bootstrap}

$\rho$ appears to be a separate postulate layered on top of the
algebraic framework.  It is not.  It is the \emph{content} of
self-modeling.

The argument has the form of a fixed-point characterization, not a
circular justification:
\begin{enumerate}[nosep]
\item Experience lives where $\rho > 0$ (by definition of $\rho$).
\item A self-modeler \emph{compresses} its experiential pattern into a
  representation --- the model, which sits at low $\rho$ (it is a
  description of experience, not experience itself).
\item The representation feeds back and \emph{becomes} experience
  (remembering, predicting, attending --- the model is experienced).
\item Return to step 1.
\end{enumerate}
The fixed point of this loop is where map equals territory.  $\rho$
defines the territory.  Self-modeling builds the map.  The fixed point
is where they coincide.  Once you have self-modeling and the algebra it forces
(Sec.~\ref{sec:why-qm}), $\rho$ should resolve to a function of the
algebra's invariants (Remark~\ref{rem:rho-h3o}).  You do not need to
postulate it separately.

Evolution provides the external anchor.  Hoffman's fitness-beats-truth
theorem~\citep{hoffman2015} shows that natural selection favors fitness
proxies over veridical world-models.  But for \emph{self}-models, the
object being modeled is the organism's own experiential state.  An
organism that captures more of its own $\rho$ is a better
self-modeler: it predicts its own states more accurately, corrects
errors faster, and maintains coherence longer.  These are survival
advantages.  Evolution does not select for truth about the
\emph{world} (Hoffman is right).  It selects for truth about the
\emph{self} ($\rho$-capture).  The organisms that survive are the ones
whose internal model best captures their own experiential measure.

The entire program in one sentence: evolution selects self-models;
self-models force the algebra; the algebra determines $\rho$; $\rho$
is what self-models model.

\subsection{The hard problem: located, not dissolved}
\label{subsec:hard-problem}

The hard problem of consciousness asks why there is ``something it is
like'' to be a self-modeling system.  The standard philosophical
framing treats this as a request for a bridge --- some additional
ingredient that converts physical processing into phenomenal
experience.

\begin{claim}
\label{claim:hard-problem}
The hard problem is \emph{self-referential}: the question is only
askable from inside the fixed point, and the asking is the
phenomenon the question asks about.
\end{claim}

``Something it is like'' is not a property that needs to be added to
the fixed point of self-modeling.  It \emph{is} what you are
describing when you describe ``hearing yourself hearing.''  The
self-reference is the experience.  Content comes from inputs (sensory
data, memories, associations).  \emph{Experiencing} comes from the
attractor closure --- the fixed point of the self-modeling
map~\citep{hofstadter1979}.  This is not dissolution (which would
declare the question confused) but \emph{location}: the question is
self-referential, and the self-reference is the phenomenon it asks
about.

\begin{remark}[Compression ratio]
The model (low $\rho$, compressed) does not feel like the experience
(high $\rho$, lived).  The explanatory gap between physical
description and phenomenal experience \emph{is} the compression
ratio.  Close the gap ($I \to H$, perfect self-knowledge) and
$\rho \to 0$: experience vanishes.  The hard problem exists precisely
because the model and the experience are at different $\rho$ levels.
The gap is the feature, not the bug.
\end{remark}
